\rhead{MN-DES: Human-Computer Interaction (HCI)}
\phantomsection
\section*{Human-Computer Interaction (HCI) (MN-DES)}
\label{minor:DES_L1}

\noindent \small{\hyperref[main:scheme]{$\leftarrow$ Back to Scheme}} \vspace{0.5cm}

\noindent \textbf{Credits:} 2 (2-0-0)

\subsection*{Course Content}
\noindent\textbf{Unit 1} \\
\textit{Human Perception and the Cognitive Model of Interaction:} HCI as an intersection of cognitive psychology, ergonomics, and software engineering; The human information processing model: sensory memory, working memory (Miller's Law: $7 \pm 2$ chunks), and long-term memory as constraints on interface design; Perception principles: Gestalt laws (proximity, similarity, continuity, closure) as the perceptual grammar of visual layout; Fitts' Law as a predictive model of pointing time: $T = a + b \log_2\!\left(\frac{2D}{W}\right)$ and its application to button sizing and menu placement; Hick's Law: decision time as a function of the number of choices; Cognitive load theory and its implications for interface complexity; Mental models: the gap between the user's conceptual model and the system's implementation model.

\vspace{0.3cm}

\noindent\textbf{Unit 2} \\
\textit{Interaction Design Paradigms and UI Architectures:} Evolution of interaction paradigms: command-line, WIMP (Windows, Icons, Menus, Pointer), touch, voice, and gesture interfaces; The Model-View-Controller (MVC) pattern reframed as a separation between data, presentation, and user intent; Event-driven programming as the computational backbone of all interactive systems: event loops, listeners, and callbacks; Affordances, signifiers, and feedback as the vocabulary of interface semantics (Norman's design principles); Information architecture: hierarchy, hub-and-spoke, and sequential navigation models; Design patterns for UI: card layouts, progressive disclosure, and contextual menus as reusable interaction solutions.

\vspace{0.3cm}

\noindent\textbf{Unit 3} \\
\textit{User Research and Requirements Engineering for Interfaces:} The user-centered design (UCD) process: discover $\rightarrow$ define $\rightarrow$ design $\rightarrow$ evaluate as an iterative loop; Qualitative research methods: contextual inquiry, semi-structured interviews, and think-aloud protocols as data collection techniques; Quantitative methods: surveys, A/B testing, and log analysis for behavioral data at scale; Personas as abstract data models representing user archetypes; User journey maps as state-machine diagrams of the user's experience over time; Task analysis: Hierarchical Task Analysis (HTA) as a tree decomposition of user goals into subtasks; Translating research findings into functional and non-functional requirements.

\vspace{0.3cm}

\noindent\textbf{Unit 4} \\
\textit{Prototyping, Wireframing, and Usability Evaluation:} The fidelity spectrum: paper sketches, wireframes, and interactive prototypes as progressively higher-cost representations; Information hierarchy and visual weight: typography scale, color contrast ratios (WCAG 2.1 AA: 4.5:1), and whitespace as layout algorithms; Usability heuristics: Nielsen's 10 principles as a checklist-based static analysis of interfaces; Usability testing: task completion rate, error rate, and time-on-task as quantitative usability metrics; System Usability Scale (SUS) as a standardized psychometric instrument; Eye-tracking and heatmaps as data visualization of visual attention patterns; Accessibility as a correctness constraint: ARIA roles, keyboard navigability, and screen reader compatibility.

\vspace{0.3cm}

\noindent\textbf{Unit 5} \\
\textit{Intelligent and Adaptive Interfaces:} Recommender systems as a form of proactive UI: collaborative filtering and content-based filtering as personalization engines; Adaptive interfaces: rule-based and ML-driven systems that modify layout, content, or interaction modality based on user context; Conversational UI: dialogue state machines as the control flow model for chatbots and voice assistants; Affective computing: inferring user emotional state from facial expression, keystroke dynamics, and physiological signals to adapt interface behavior; Dark patterns as adversarial HCI: taxonomizing manipulative design and their detection as a classification problem; Ethical dimensions of persuasive technology: nudging, attention economics, and the responsibility of interface engineers.

\newpage
\rhead{Curriculum Schema}